\IfFileExists{stacks-project.cls}{%
\documentclass{stacks-project}
}{%
\documentclass{amsart}
}

% For dealing with references we use the comment environment
\usepackage{verbatim}
\newenvironment{reference}{\comment}{\endcomment}
%\newenvironment{reference}{}{}
\newenvironment{slogan}{\comment}{\endcomment}
\newenvironment{history}{\comment}{\endcomment}

% For commutative diagrams we use Xy-pic
\usepackage[all]{xy}

% We use 2cell for 2-commutative diagrams.
\xyoption{2cell}
\UseAllTwocells

% We use multicol for the list of chapters between chapters
\usepackage{multicol}

% This is generally recommended for better output
\usepackage{lmodern}
\usepackage[T1]{fontenc}

% For cross-file-references
\usepackage{xr-hyper}

% Package for hypertext links:
\usepackage{hyperref}

% For any local file, say "hello.tex" you want to link to please
% use \externaldocument[hello-]{hello}
\externaldocument[introduction-]{introduction}
\externaldocument[conventions-]{conventions}
\externaldocument[sets-]{sets}
\externaldocument[categories-]{categories}
\externaldocument[topology-]{topology}
\externaldocument[sheaves-]{sheaves}
\externaldocument[sites-]{sites}
\externaldocument[stacks-]{stacks}
\externaldocument[fields-]{fields}
\externaldocument[algebra-]{algebra}
\externaldocument[brauer-]{brauer}
\externaldocument[homology-]{homology}
\externaldocument[derived-]{derived}
\externaldocument[simplicial-]{simplicial}
\externaldocument[more-algebra-]{more-algebra}
\externaldocument[smoothing-]{smoothing}
\externaldocument[modules-]{modules}
\externaldocument[sites-modules-]{sites-modules}
\externaldocument[injectives-]{injectives}
\externaldocument[cohomology-]{cohomology}
\externaldocument[sites-cohomology-]{sites-cohomology}
\externaldocument[dga-]{dga}
\externaldocument[dpa-]{dpa}
\externaldocument[sdga-]{sdga}
\externaldocument[hypercovering-]{hypercovering}
\externaldocument[schemes-]{schemes}
\externaldocument[constructions-]{constructions}
\externaldocument[properties-]{properties}
\externaldocument[morphisms-]{morphisms}
\externaldocument[coherent-]{coherent}
\externaldocument[divisors-]{divisors}
\externaldocument[limits-]{limits}
\externaldocument[varieties-]{varieties}
\externaldocument[topologies-]{topologies}
\externaldocument[descent-]{descent}
\externaldocument[perfect-]{perfect}
\externaldocument[more-morphisms-]{more-morphisms}
\externaldocument[flat-]{flat}
\externaldocument[groupoids-]{groupoids}
\externaldocument[more-groupoids-]{more-groupoids}
\externaldocument[etale-]{etale}
\externaldocument[chow-]{chow}
\externaldocument[intersection-]{intersection}
\externaldocument[pic-]{pic}
\externaldocument[weil-]{weil}
\externaldocument[adequate-]{adequate}
\externaldocument[dualizing-]{dualizing}
\externaldocument[duality-]{duality}
\externaldocument[discriminant-]{discriminant}
\externaldocument[derham-]{derham}
\externaldocument[local-cohomology-]{local-cohomology}
\externaldocument[algebraization-]{algebraization}
\externaldocument[curves-]{curves}
\externaldocument[resolve-]{resolve}
\externaldocument[models-]{models}
\externaldocument[functors-]{functors}
\externaldocument[equiv-]{equiv}
\externaldocument[pione-]{pione}
\externaldocument[etale-cohomology-]{etale-cohomology}
\externaldocument[proetale-]{proetale}
\externaldocument[relative-cycles-]{relative-cycles}
\externaldocument[more-etale-]{more-etale}
\externaldocument[trace-]{trace}
\externaldocument[crystalline-]{crystalline}
\externaldocument[spaces-]{spaces}
\externaldocument[spaces-properties-]{spaces-properties}
\externaldocument[spaces-morphisms-]{spaces-morphisms}
\externaldocument[decent-spaces-]{decent-spaces}
\externaldocument[spaces-cohomology-]{spaces-cohomology}
\externaldocument[spaces-limits-]{spaces-limits}
\externaldocument[spaces-divisors-]{spaces-divisors}
\externaldocument[spaces-over-fields-]{spaces-over-fields}
\externaldocument[spaces-topologies-]{spaces-topologies}
\externaldocument[spaces-descent-]{spaces-descent}
\externaldocument[spaces-perfect-]{spaces-perfect}
\externaldocument[spaces-more-morphisms-]{spaces-more-morphisms}
\externaldocument[spaces-flat-]{spaces-flat}
\externaldocument[spaces-groupoids-]{spaces-groupoids}
\externaldocument[spaces-more-groupoids-]{spaces-more-groupoids}
\externaldocument[bootstrap-]{bootstrap}
\externaldocument[spaces-pushouts-]{spaces-pushouts}
\externaldocument[spaces-chow-]{spaces-chow}
\externaldocument[groupoids-quotients-]{groupoids-quotients}
\externaldocument[spaces-more-cohomology-]{spaces-more-cohomology}
\externaldocument[spaces-simplicial-]{spaces-simplicial}
\externaldocument[spaces-duality-]{spaces-duality}
\externaldocument[formal-spaces-]{formal-spaces}
\externaldocument[restricted-]{restricted}
\externaldocument[spaces-resolve-]{spaces-resolve}
\externaldocument[formal-defos-]{formal-defos}
\externaldocument[defos-]{defos}
\externaldocument[cotangent-]{cotangent}
\externaldocument[examples-defos-]{examples-defos}
\externaldocument[algebraic-]{algebraic}
\externaldocument[examples-stacks-]{examples-stacks}
\externaldocument[stacks-sheaves-]{stacks-sheaves}
\externaldocument[criteria-]{criteria}
\externaldocument[artin-]{artin}
\externaldocument[quot-]{quot}
\externaldocument[stacks-properties-]{stacks-properties}
\externaldocument[stacks-morphisms-]{stacks-morphisms}
\externaldocument[stacks-limits-]{stacks-limits}
\externaldocument[stacks-cohomology-]{stacks-cohomology}
\externaldocument[stacks-perfect-]{stacks-perfect}
\externaldocument[stacks-introduction-]{stacks-introduction}
\externaldocument[stacks-more-morphisms-]{stacks-more-morphisms}
\externaldocument[stacks-geometry-]{stacks-geometry}
\externaldocument[moduli-]{moduli}
\externaldocument[moduli-curves-]{moduli-curves}
\externaldocument[examples-]{examples}
\externaldocument[exercises-]{exercises}
\externaldocument[guide-]{guide}
\externaldocument[desirables-]{desirables}
\externaldocument[coding-]{coding}
\externaldocument[obsolete-]{obsolete}
\externaldocument[fdl-]{fdl}
\externaldocument[index-]{index}

% Theorem environments.
%
\theoremstyle{plain}
\newtheorem{theorem}[subsection]{Theorem}
\newtheorem{proposition}[subsection]{Proposition}
\newtheorem{lemma}[subsection]{Lemma}

\theoremstyle{definition}
\newtheorem{definition}[subsection]{Definition}
\newtheorem{example}[subsection]{Example}
\newtheorem{exercise}[subsection]{Exercise}
\newtheorem{situation}[subsection]{Situation}

\theoremstyle{remark}
\newtheorem{remark}[subsection]{Remark}
\newtheorem{remarks}[subsection]{Remarks}

\numberwithin{equation}{subsection}

% Macros
%
\def\lim{\mathop{\mathrm{lim}}\nolimits}
\def\colim{\mathop{\mathrm{colim}}\nolimits}
\def\Spec{\mathop{\mathrm{Spec}}}
\def\Hom{\mathop{\mathrm{Hom}}\nolimits}
\def\Ext{\mathop{\mathrm{Ext}}\nolimits}
\def\SheafHom{\mathop{\mathcal{H}\!\mathit{om}}\nolimits}
\def\SheafExt{\mathop{\mathcal{E}\!\mathit{xt}}\nolimits}
\def\Sch{\mathit{Sch}}
\def\Mor{\mathop{\mathrm{Mor}}\nolimits}
\def\Ob{\mathop{\mathrm{Ob}}\nolimits}
\def\Sh{\mathop{\mathit{Sh}}\nolimits}
\def\NL{\mathop{N\!L}\nolimits}
\def\CH{\mathop{\mathrm{CH}}\nolimits}
\def\proetale{{pro\text{-}\acute{e}tale}}
\def\etale{{\acute{e}tale}}
\def\QCoh{\mathit{QCoh}}
\def\Ker{\mathop{\mathrm{Ker}}}
\def\Im{\mathop{\mathrm{Im}}}
\def\Coker{\mathop{\mathrm{Coker}}}
\def\Coim{\mathop{\mathrm{Coim}}}

% Boxtimes
%
\DeclareMathSymbol{\boxtimes}{\mathbin}{AMSa}{"02}

%
% Macros for moduli stacks/spaces
%
\def\QCohstack{\mathcal{QC}\!\mathit{oh}}
\def\Cohstack{\mathcal{C}\!\mathit{oh}}
\def\Spacesstack{\mathcal{S}\!\mathit{paces}}
\def\Quotfunctor{\mathrm{Quot}}
\def\Hilbfunctor{\mathrm{Hilb}}
\def\Curvesstack{\mathcal{C}\!\mathit{urves}}
\def\Polarizedstack{\mathcal{P}\!\mathit{olarized}}
\def\Complexesstack{\mathcal{C}\!\mathit{omplexes}}
% \Pic is the operator that assigns to X its picard group, usage \Pic(X)
% \Picardstack_{X/B} denotes the Picard stack of X over B
% \Picardfunctor_{X/B} denotes the Picard functor of X over B
\def\Pic{\mathop{\mathrm{Pic}}\nolimits}
\def\Picardstack{\mathcal{P}\!\mathit{ic}}
\def\Picardfunctor{\mathrm{Pic}}
\def\Deformationcategory{\mathcal{D}\!\mathit{ef}}
\usepackage{cleveref}
\crefname{originalreference}{original reference}{original references}
\Crefname{originalreference}{Original reference}{Original references}
\makeatletter
\newlabel{item-axiom-1}{{0AND}{source}{}{sourcepointer-1}{}}
\newlabel{item-axiom-1@cref}{{[originalreference][0][]0AND}{[1][0][]source}{}{}{}}
\makeatother
\crefname{originalreference}{original reference}{original references}
\Crefname{originalreference}{Original reference}{Original references}
\makeatletter
\newlabel{item-axiom-2}{{0ANE}{source}{}{sourcepointer-2}{}}
\newlabel{item-axiom-2@cref}{{[originalreference][0][]0ANE}{[1][0][]source}{}{}{}}
\makeatother
\crefname{originalreference}{original reference}{original references}
\Crefname{originalreference}{Original reference}{Original references}
\makeatletter
\newlabel{item-axiom-3}{{0ANF}{source}{}{sourcepointer-3}{}}
\newlabel{item-axiom-3@cref}{{[originalreference][0][]0ANF}{[1][0][]source}{}{}{}}
\makeatother
\crefname{originallemma}{original lemma}{original lemmas}
\Crefname{originallemma}{Original lemma}{Original lemmas}
\makeatletter
\newlabel{lemma-etale}{{0AN8}{source}{}{sourcepointer-4}{}}
\newlabel{lemma-etale@cref}{{[originallemma][0][]0AN8}{[1][0][]source}{}{}{}}
\makeatother
\crefname{originallemma}{original lemma}{original lemmas}
\Crefname{originallemma}{Original lemma}{Original lemmas}
\makeatletter
\newlabel{lemma-etale-surjective}{{0AN9}{source}{}{sourcepointer-5}{}}
\newlabel{lemma-etale-surjective@cref}{{[originallemma][0][]0AN9}{[1][0][]source}{}{}{}}
\makeatother
\crefname{originallemma}{original lemma}{original lemmas}
\Crefname{originallemma}{Original lemma}{Original lemmas}
\makeatletter
\newlabel{lemma-fibre-products-general}{{0AJ2}{source}{}{sourcepointer-6}{}}
\newlabel{lemma-fibre-products-general@cref}{{[originallemma][0][]0AJ2}{[1][0][]source}{}{}{}}
\makeatother
\crefname{originalremark}{original remark}{original remarks}
\Crefname{originalremark}{Original remark}{Original remarks}
\makeatletter
\newlabel{remark-set-theoretic}{{0AIS}{source}{}{sourcepointer-7}{}}
\newlabel{remark-set-theoretic@cref}{{[originalremark][0][]0AIS}{[1][0][]source}{}{}{}}
\makeatother
\title{Three equivalent formulations of a local property of morphisms of complete weakly admissible rings}
\author{The Stacks Project Authors}
\date{stacks-project:formal-spaces.tex:local-snapshot-20261012; GNU FDL 1.2 or later}
\begin{document}
\maketitle
\noindent\textit{Editorial note.} All delivered mathematical text is the author's own, copied from the pinned Stacks Project chapter file formal-spaces.tex; no statement or proof step was rewritten, reordered or removed.\par\medskip
\noindent\textit{Editorial note.} The proof displays two large commutative diagrams and one small one; all lie inside the counted span and are reproduced verbatim.\par\medskip
\noindent\textit{Editorial note.} The axioms are referred to by label; those references lie outside the counted span and are carried as navigation stubs rather than reproduced, and the situation governing them is recorded as a standard prerequisite.\par\medskip
\noindent\textit{Editorial note.} The proof contains the author's parenthetical noting that applying the second axiom with identities is permissible because identities are faithfully flat ring maps; it is reproduced verbatim.\par\medskip
\noindent\textit{Editorial note.} The counted spans carry 8 navigation stubs; no bibliography citation occurs in the spans.\par\medskip
\noindent\textit{Editorial note.} The preamble span points at the chapter's own preamble.tex shipped with the frozen source, used verbatim as staged.\par\medskip
\noindent\textit{Allowed standard prerequisites (editorial).}\begin{itemize}
\item The locality axioms for a local property of morphisms of complete weakly admissible rings. Same chapter, the situation governing these axioms, referred to at proof lines 5155-5156, 5160-5166 and 5171-5175, with the individual items item-axiom-1, item-axiom-2 and item-axiom-3. Supply the three moves by which the property is transported around the diagram.
\item The definition of a formal algebraic space through coverings, and fibre products of formal algebraic spaces. Same chapter, Definition definition-formal-algebraic-space used throughout at proof lines 5119-5126, and Lemma lemma-fibre-products-general at proof lines 5133-5135. Provide the coverings whose members are affine and the statement that the fibre products appearing in the chase are formal algebraic spaces.
\item Morphisms between affine formal algebraic spaces coming from completions of etale ring maps, and surjective such morphisms giving faithfully flat etale ring maps. Same chapter, Lemma lemma-etale at proof lines 5150-5153 and Lemma lemma-etale-surjective at proof lines 5177-5179. Identify the ring maps to which the axioms are applied and recognise the final covering as a faithfully flat etale extension.
\end{itemize}\medskip
\ifdefined\hypertarget\hypertarget{sourcepointer-1}{}\fi
\noindent\textit{Source reference.} Original reference 0AND: sources/stacks-formal-spaces/pinned-original-source.tex, source lines 5086--5086. Pointer outside the counted span, invoked once for context; the load-bearing argument of the counted statement is the author's proof reproduced in full in the proof span.\par\medskip
\ifdefined\hypertarget\hypertarget{sourcepointer-2}{}\fi
\noindent\textit{Source reference.} Original reference 0ANE: sources/stacks-formal-spaces/pinned-original-source.tex, source lines 5090--5090. Pointer outside the counted span, invoked once for context; the load-bearing argument of the counted statement is the author's proof reproduced in full in the proof span.\par\medskip
\ifdefined\hypertarget\hypertarget{sourcepointer-3}{}\fi
\noindent\textit{Source reference.} Original reference 0ANF: sources/stacks-formal-spaces/pinned-original-source.tex, source lines 5095--5095. Pointer outside the counted span, invoked once for context; the load-bearing argument of the counted statement is the author's proof reproduced in full in the proof span.\par\medskip
\ifdefined\hypertarget\hypertarget{sourcepointer-4}{}\fi
\noindent\textit{Source reference.} Original lemma 0AN8: sources/stacks-algebra/pinned-original-source.tex, source lines 39507--39507. Pointer outside the counted span, invoked once for context; the load-bearing argument of the counted statement is the author's proof reproduced in full in the proof span.\par\medskip
\ifdefined\hypertarget\hypertarget{sourcepointer-5}{}\fi
\noindent\textit{Source reference.} Original lemma 0AN9: sources/stacks-formal-spaces/pinned-original-source.tex, source lines 4492--4492. Pointer outside the counted span, invoked once for context; the load-bearing argument of the counted statement is the author's proof reproduced in full in the proof span.\par\medskip
\ifdefined\hypertarget\hypertarget{sourcepointer-6}{}\fi
\noindent\textit{Source reference.} Original lemma 0AJ2: sources/stacks-formal-spaces/pinned-original-source.tex, source lines 3488--3488. Pointer outside the counted span, invoked once for context; the load-bearing argument of the counted statement is the author's proof reproduced in full in the proof span.\par\medskip
\ifdefined\hypertarget\hypertarget{sourcepointer-7}{}\fi
\noindent\textit{Source reference.} Original remark 0AIS: sources/stacks-formal-spaces/pinned-original-source.tex, source lines 2769--2769. Pointer outside the counted span, invoked once for context; the load-bearing argument of the counted statement is the author's proof reproduced in full in the proof span.\par\medskip
\begin{definition}
\label{definition-formal-algebraic-space}
Let $S$ be a scheme. We say a sheaf $X$ on $(\Sch/S)_{fppf}$ is a
{\it formal algebraic space} if there exist a family of maps
$\{X_i \to X\}_{i \in I}$ of sheaves such that
\begin{enumerate}
\item $X_i$ is an affine formal algebraic space,
\item $X_i \to X$ is representable by algebraic spaces and \'etale,
\item $\coprod X_i \to X$ is surjective as a map of sheaves
\end{enumerate}
and $X$ satisfies a set theoretic condition
(see Remark \ref{remark-set-theoretic}). A
{\it morphism of formal algebraic spaces}
over $S$ is a map of sheaves.
\end{definition}

\begin{lemma}
\label{lemma-property-defines-property-morphisms}
Let $S$ be a scheme. Let $f : X \to Y$ be a morphism of
locally countably indexed formal algebraic spaces over $S$.
Let $P$ be a local property of morphisms of $\textit{WAdm}^{count}$.
The following are equivalent
\begin{enumerate}
\item for every commutative diagram
$$
\xymatrix{
U \ar[d] \ar[r] & V \ar[d] \\
X \ar[r] & Y
}
$$
with $U$ and $V$ affine formal algebraic spaces, $U \to X$ and $V \to Y$
representable by algebraic spaces and \'etale, the morphism $U \to V$
corresponds to a morphism of $\textit{WAdm}^{count}$ with property $P$,
\item there exists a covering $\{Y_j \to Y\}$ as in
Definition \ref{definition-formal-algebraic-space} and for each $j$
a covering $\{X_{ji} \to Y_j \times_Y X\}$ as in
Definition \ref{definition-formal-algebraic-space}
such that each $X_{ji} \to Y_j$  corresponds
to a morphism of $\textit{WAdm}^{count}$ with property $P$, and
\item there exist a covering $\{X_i \to X\}$ as in
Definition \ref{definition-formal-algebraic-space}
and for each $i$ a factorization $X_i \to Y_i \to Y$ where $Y_i$
is an affine formal algebraic space, $Y_i \to Y$ is representable
by algebraic spaces and \'etale, and $X_i \to Y_i$ corresponds
to a morphism of $\textit{WAdm}^{count}$ with property $P$.
\end{enumerate}
\end{lemma}

\begin{proof}
It is clear that (1) implies (2) and that (2) implies (3).
Assume $\{X_i \to X\}$ and $X_i \to Y_i \to Y$ as in (3)
and let a diagram as in (1) be given.
Since $Y_i \times_Y V$ is a formal algebraic space
(Lemma \ref{lemma-fibre-products-general}) we may pick
coverings $\{Y_{ij} \to Y_i \times_Y V\}$ as in
Definition \ref{definition-formal-algebraic-space}.
For each $(i, j)$ we may similarly choose coverings
$\{X_{ijk} \to Y_{ij} \times_{Y_i} X_i \times_X U\}$
as in Definition \ref{definition-formal-algebraic-space}.
Since $U$ is quasi-compact we can choose
$(i_1, j_1, k_1), \ldots, (i_n, j_n, k_n)$ such that
$$
X_{i_1 j_1 k_1} \amalg \ldots \amalg X_{i_n j_n k_n} \longrightarrow U
$$
is surjective. For $s = 1, \ldots, n$ consider the commutative diagram
$$
\xymatrix{
& & & X_{i_s j_s k_s} \ar[ld] \ar[d] \ar[rd] \\
X \ar[d] & X_{i_s} \ar[l] \ar[d] &
X_{i_s} \times_X U \ar[l] \ar[d] & Y_{i_s j_s} \ar[ld] \ar[rd] &
X_{i_s} \times_X U \ar[d] \ar[r] &
U \ar[d] \ar[r] & X \ar[d] \\
Y & Y_{i_s} \ar[l] &
Y_{i_s} \times_Y V \ar[l] & &
Y_{i_s} \times_Y V \ar[r] &
V \ar[r] & Y
}
$$
Let us say that $P$ holds for a morphism of countably indexed
affine formal algebraic spaces if it holds for the corresponding
morphism of $\textit{WAdm}^{count}$. Observe that the maps
$X_{i_s j_s k_s} \to X_{i_s}$, $Y_{i_s j_s} \to Y_{i_s}$
are given by completions of \'etale ring maps, see Lemma \ref{lemma-etale}.
Hence we see that $P(X_{i_s} \to Y_{i_s})$ implies
$P(X_{i_s j_s k_s} \to Y_{i_s j_s})$ by axiom (\ref{item-axiom-1}).
Observe that the maps $Y_{i_s j_s} \to V$ are given by completions of
\'etale rings maps (same lemma as before).
By axiom (\ref{item-axiom-2}) applied to the diagram
$$
\xymatrix{
X_{i_s j_s k_s} \ar@{=}[r] \ar[d] & X_{i_s j_s k_s} \ar[d] \\
Y_{i_s j_s} \ar[r] & V
}
$$
(this is permissible as identities are faithfully flat ring maps)
we conclude that $P(X_{i_s j_s k_s} \to V)$ holds.
By axiom (\ref{item-axiom-3}) we find that
$P(\coprod_{s = 1, \ldots, n} X_{i_s j_s k_s} \to V)$ holds.
Since the morphism $\coprod X_{i_s j_s k_s} \to U$ is surjective
by construction, the corresponding morphism of $\textit{WAdm}^{count}$
is the completion of a faithfully flat \'etale ring map, see
Lemma \ref{lemma-etale-surjective}.
One more application of axiom (\ref{item-axiom-2})
(with $B' = B$) implies that $P(U \to V)$ is true as desired.
\end{proof}

\end{document}
